\newpage
\section{Introduction}

\subsection{Motivations}

Occupational choice is among the most consequential decisions an individual makes, and it is typically made with limited direct knowledge of the work itself. In the taxonomy of Gati, Krausz and Osipow\cite{r1}, whose structure was recently confirmed across thirteen countries\cite{r2}, insufficient information about occupations constitutes a principal category of decision difficulty.

\subsection{Objectives}

The overall objective is to enable a person facing an occupational decision to experience the work of an occupation before committing to one. It is pursued through two specific objectives:

\begin{enumerate}[leftmargin=1.5em]
  \item To deliver occupational tasks that remain faithful to the occupations they represent.
  \item To assess users' responses through automated judgment whose agreement with domain experts is established.
\end{enumerate}

\subsection{Scope}

The system is complete in function and bounded in coverage to the information technology field.

\subsection{Significance of the Project}

\subsubsection{Significance from the Practical Perspectives}

The practical contribution is a freely accessible means of acquiring occupational information experientially before commitment.

\subsubsection{Significance from the Scientific Perspectives}

The scientific contribution lies in grounded generation of educational content and validated automated assessment.

\subsection{Report Structure}

Chapter 2 establishes the background, Chapter 3 reviews related work, and Chapter 4 presents the analysis and design.
